% GENERATED FILE. Edit canonical lessons, then run npm run generate:books.
% canonical-source-hash: fnv1a64:d4ff3ae5ce04a092
% canonical-lessons: HI-C238-kacca, HI-C238-gunguna, HI-C238-rukha, HI-C238-cipcipa, HI-C238-khurdura

\chapter{Raw, Lukewarm, Dry, Sticky, Rough}
\label{ch:hi-raw-lukewarm-dry-sticky-rough}
\hlchaptermodality{238}

\begin{chapteropening}
\textbf{By the end of this chapter:} \emph{I can say raw, lukewarm, dry, sticky and rough.}

Complete the last lesson of chapter 238: I can say raw, lukewarm, dry, sticky and rough.
\end{chapteropening}

\section[\texorpdfstring{\hi{कच्चा} (kaccā)}{कच्चा (kaccā)}]{\hi{कच्चा} (kaccā) — raw, unripe}

\label{lesson:HI-C238-kacca}



\begin{quote}
Before the new one: say the Hindi for shallow, then the Hindi for stale.
\end{quote}

\subsection*{You'll want to know: \hi{कच्चा}}
\textbf{\hi{कच्चा}} — \emph{kaccā} — \textquotedblleft{}raw, unripe\textquotedblright{}.

\begin{grammarlens}[title={What you've built}]
One more word for this chapter.
\end{grammarlens}

\subsection*{Your turn}
\begin{itemize}
\raggedright
  \item \emph{Say it:} \emph{kaccā}
  \item \emph{Say it:} \emph{kaccā}, once more
  \item \emph{Say it:} say \emph{kaccā}
  \item \emph{Recall:} say the Hindi for shallow, then the Hindi for stale, then say \emph{kaccā} again
\end{itemize}

\begin{tcolorbox}[breakable,colback=teal!4,colframe=teal!35!black,title={Before you move on}]
What does \hi{कच्चा} mean? (\textquotedblleft{}Raw, unripe\textquotedblright{}.) Say it once more.
\end{tcolorbox}
\section[\texorpdfstring{\hi{गुनगुना} (gungunā)}{गुनगुना (gungunā)}]{\hi{गुनगुना} (gungunā) — lukewarm}

\label{lesson:HI-C238-gunguna}



\begin{quote}
Before the new one: say the Hindi for stale, then the Hindi for raw.
\end{quote}

\subsection*{You'll want to know: \hi{गुनगुना}}
\textbf{\hi{गुनगुना}} — \emph{gungunā} — \textquotedblleft{}lukewarm\textquotedblright{}.

\begin{grammarlens}[title={What you've built}]
One more word for this chapter.
\end{grammarlens}

\subsection*{Your turn}
\begin{itemize}
\raggedright
  \item \emph{Say it:} \emph{gungunā}
  \item \emph{Say it:} \emph{gungunā}, once more
  \item \emph{Say it:} say \emph{gungunā}
  \item \emph{Recall:} say the Hindi for stale, then the Hindi for raw, then say \emph{gungunā} again
\end{itemize}

\begin{tcolorbox}[breakable,colback=teal!4,colframe=teal!35!black,title={Before you move on}]
What does \hi{गुनगुना} mean? (\textquotedblleft{}Lukewarm\textquotedblright{}.) Say it once more.
\end{tcolorbox}
\section[\texorpdfstring{\hi{रूखा} (rūkhā)}{रूखा (rūkhā)}]{\hi{रूखा} (rūkhā) — dry, rough}

\label{lesson:HI-C238-rukha}



\begin{quote}
Before the new one: say the Hindi for raw, then the Hindi for lukewarm.
\end{quote}

\subsection*{You'll want to know: \hi{रूखा}}
\textbf{\hi{रूखा}} — \emph{rūkhā} — \textquotedblleft{}dry, rough\textquotedblright{}.

\begin{grammarlens}[title={What you've built}]
One more word for this chapter.
\end{grammarlens}

\subsection*{Your turn}
\begin{itemize}
\raggedright
  \item \emph{Say it:} \emph{rūkhā}
  \item \emph{Say it:} \emph{rūkhā}, once more
  \item \emph{Say it:} say \emph{rūkhā}
  \item \emph{Recall:} say the Hindi for raw, then the Hindi for lukewarm, then say \emph{rūkhā} again
\end{itemize}

\begin{tcolorbox}[breakable,colback=teal!4,colframe=teal!35!black,title={Before you move on}]
What does \hi{रूखा} mean? (\textquotedblleft{}Dry, rough\textquotedblright{}.) Say it once more.
\end{tcolorbox}
\section[\texorpdfstring{\hi{चिपचिपा} (cipcipā)}{चिपचिपा (cipcipā)}]{\hi{चिपचिपा} (cipcipā) — sticky}

\label{lesson:HI-C238-cipcipa}



\begin{quote}
Before the new one: say the Hindi for lukewarm, then the Hindi for dry.
\end{quote}

\subsection*{You'll want to know: \hi{चिपचिपा}}
\textbf{\hi{चिपचिपा}} — \emph{cipcipā} — \textquotedblleft{}sticky\textquotedblright{}.

\begin{grammarlens}[title={What you've built}]
One more word for this chapter.
\end{grammarlens}

\subsection*{Your turn}
\begin{itemize}
\raggedright
  \item \emph{Say it:} \emph{cipcipā}
  \item \emph{Say it:} \emph{cipcipā}, once more
  \item \emph{Say it:} say \emph{cipcipā}
  \item \emph{Recall:} say the Hindi for lukewarm, then the Hindi for dry, then say \emph{cipcipā} again
\end{itemize}

\begin{tcolorbox}[breakable,colback=teal!4,colframe=teal!35!black,title={Before you move on}]
What does \hi{चिपचिपा} mean? (\textquotedblleft{}Sticky\textquotedblright{}.) Say it once more.
\end{tcolorbox}
\section[\texorpdfstring{\hi{खुरदुरा} (khurdurā)}{खुरदुरा (khurdurā)}]{\hi{खुरदुरा} (khurdurā) — rough (to the touch)}

\label{lesson:HI-C238-khurdura}



\begin{quote}
Before the new one: say the Hindi for dry, then the Hindi for sticky.
\end{quote}

\subsection*{You'll want to know: \hi{खुरदुरा}}
\textbf{\hi{खुरदुरा}} — \emph{khurdurā} — \textquotedblleft{}rough (to the touch)\textquotedblright{}.

\begin{grammarlens}[title={What you've built}]
One more word for this chapter.
\end{grammarlens}

\subsection*{Your turn}
\begin{itemize}
\raggedright
  \item \emph{Say it:} \emph{khurdurā}
  \item \emph{Say it:} \emph{khurdurā}, once more
  \item \emph{Say it:} say \emph{khurdurā}
  \item \emph{Recall:} say the Hindi for dry, then the Hindi for sticky, then say \emph{khurdurā} again
\end{itemize}

\begin{tcolorbox}[breakable,colback=teal!4,colframe=teal!35!black,title={Before you move on}]
What does \hi{खुरदुरा} mean? (\textquotedblleft{}Rough (to the touch)\textquotedblright{}.) Say it once more.
\end{tcolorbox}
